\section{Session headers and framing}\label{sec:session_headers}

Every session-layer message begins with a fixed-size header of 24 bytes,
followed by a type-specific payload. Multi-byte integer fields are encoded in
little-endian byte order and are positioned to favour natural alignment on
four- and eight-byte boundaries. The wire layouts below are given in DSDL
notation.
Fields marked \emph{incompatibility} shall be transmitted as
zero, and any non-zero value on reception indicates a protocol revision the
receiver does not understand: the receiver shall discard the message.

\subsection{Header types}

The first byte of every header is the header type. Types are allocated as
follows; values 10 through 255 are reserved.

{\footnotesize
\begin{CyphalSimpleTable}[wide]{Session-layer header types}{|r l l X|}
    Value & Name             & Direction                         & Purpose                                        \\
    0     & \texttt{msg\_be}    & publisher $\to$ subscriber(s), multicast or unicast & Best-effort message publication.                \\
    1     & \texttt{msg\_rel}   & publisher $\to$ subscriber(s), multicast or unicast & Reliable message publication, acknowledged.     \\
    2     & \texttt{msg\_ack}   & subscriber $\to$ publisher, unicast              & Positive acknowledgement of a reliable message. \\
    3     & \texttt{msg\_nack}  & subscriber $\to$ publisher, unicast              & Negative acknowledgement; no matching subscriber.\\
    4     & \texttt{rsp\_be}    & responder $\to$ publisher, unicast               & Best-effort response to a received message.     \\
    5     & \texttt{rsp\_rel}   & responder $\to$ publisher, unicast               & Reliable response, acknowledged.                \\
    6     & \texttt{rsp\_ack}   & publisher $\to$ responder, unicast               & Positive acknowledgement of a reliable response.\\
    7     & \texttt{rsp\_nack}  & publisher $\to$ responder, unicast               & Negative acknowledgement; application no longer listening. \\
    8     & \texttt{gossip}     & any, multicast or unicast                        & Topic allocation CRDT gossip.                   \\
    9     & \texttt{scout}      & any, broadcast                                   & Pattern service-discovery query.                \\
\end{CyphalSimpleTable}
}

\subsection{Message headers}\label{sec:session_header_msg}

The \texttt{msg\_be} and \texttt{msg\_rel} headers share the same layout. Each
message carries not only the application payload but also the sender's
current CRDT state for the topic, so that the message doubles as an inline
gossip for its own topic (section~\ref{sec:session_gossip}).

\begin{minted}{python}
uint8   type             # 0 for msg_be, 1 for msg_rel.
void8
uint8   incompatibility
int8    topic_log_age    # floor(log2(topic_age_seconds)) if topic_age_seconds > 0, else -1.
uint32  evictions        # For pinned topics, 0xFFFFFFFF - subject_id.
uint64  topic_hash       # Subject allocation collision detection and immediate consensus updates.
uint64  tag              # Random-initialized, wraps around. Used for acknowledgement
                         # correlation and, optionally, for ordering recovery per remote.
# Payload follows.
\end{minted}

\subsection{Message acknowledgement headers}\label{sec:session_header_msg_ack}

The \texttt{msg\_ack} and \texttt{msg\_nack} headers are sent in response to a
\texttt{msg\_rel}; they carry no payload. The priority level of the
acknowledgement shall match the priority level of the original message.

\begin{minted}{python}
uint8   type             # 2 for msg_ack, 3 for msg_nack.
void24
uint32  incompatibility
uint64  topic_hash       # From the acknowledged message.
uint64  tag              # From the acknowledged message.
\end{minted}

\subsection{Response headers}\label{sec:session_header_rsp}

The \texttt{rsp\_be} and \texttt{rsp\_rel} headers carry a response to a
message previously received by the responder. The \texttt{rsp\_ack} and
\texttt{rsp\_nack} headers acknowledge a reliable response and carry no
payload. All four share the same layout.

\begin{minted}{python}
uint8   type             # 4 for rsp_be, 5 for rsp_rel, 6 for rsp_ack, 7 for rsp_nack.
uint8   tag              # Chosen by the responder arbitrarily for acknowledgement correlation.
uint48  seqno            # Starts at zero for the first response to a given message and
                         # increments by one per successive response, enabling streaming.
uint64  topic_hash       # Topic hash of the published message this response pertains to.
uint64  message_tag      # Tag of the published message this response pertains to.
# Payload follows, except for rsp_ack and rsp_nack.
\end{minted}

\subsection{Gossip header}\label{sec:session_header_gossip}

The \texttt{gossip} header carries one CRDT record
(section~\ref{sec:session_crdt}) together with the topic's normalized
resolved name. The message has no additional payload.

\begin{minted}{python}
uint8   type             # 8.
void16
int8    topic_log_age    # floor(log2(topic_age_seconds)) if topic_age_seconds > 0, else -1.
uint32  incompatibility
uint64  topic_hash
uint32  topic_evictions  # For pinned topics, 0xFFFFFFFF - subject_id.
void24                   # Reserved, may extend the eviction counter in a future revision.
utf8[<=200] topic_name   # Length-prefixed (one byte prefix); normalized resolved name.
# Total wire size is 24 bytes plus the length of the topic name.
\end{minted}

\subsection{Scout header}\label{sec:session_header_scout}

The \texttt{scout} header carries a pattern string
(section~\ref{sec:session_names}). The receiver matches the pattern against
its local topics and, for each match, replies with a \texttt{gossip} message
(section~\ref{sec:session_scout}).

\begin{minted}{python}
uint8   type             # 9.
void24
uint32  incompatibility
uint64  incompatibility1
void56
utf8[<=200] pattern      # Length-prefixed (one byte prefix); pattern applied to normalized names.
# Total wire size is 24 bytes plus the length of the pattern.
\end{minted}

\subsection{Message tags}\label{sec:session_tags}

The 64-bit \texttt{tag} field carried in \texttt{msg\_be} and \texttt{msg\_rel}
is a per-publisher counter that is randomly initialized when a publisher is
created and is incremented by one for each message published on the topic.
The random initialization ensures that a node rebooting mid-operation does not
accidentally reuse tags that were emitted before the reboot; the counter is
permitted to wrap around.

The PRNG used for tag initialization shall have distinct state across quick
reboots. Nodes without a true random-number generator should seed a pseudorandom
generator by hashing the node's unique identifier with whatever entropy the
platform makes available -- for example, analog-to-digital converter noise,
real-time clock drift, or values stored in battery-backed memory.

Tags are not used for duplicate suppression on best-effort messages, which rely
on the transport layer's deduplication instead. They are used for duplicate
suppression and acknowledgement correlation on reliable messages
(section~\ref{sec:session_reliable}).

\section{Best-effort message exchange}\label{sec:session_best_effort}

A best-effort publication is a single \texttt{msg\_be} transfer. The publisher
performs the following steps atomically:

\begin{enumerate}
    \item Assign the next tag from the publisher's sequence.
    \item Serialize a \texttt{msg\_be} header whose inline CRDT fields carry
    the topic's current $\mathit{lage}$, eviction counter, and hash.
    \item Append the application payload.
    \item Submit the message to the transport for delivery on the topic's
    current subject-ID.
\end{enumerate}

The publisher retains no state after the submission returns.

On reception of a \texttt{msg\_be} transfer, a subscriber shall:

\begin{enumerate}
    \item Validate that the header's \texttt{(topic\_hash, evictions)} pair
    maps (via section~\ref{sec:session_subject_id_mapping}) to the subject-ID
    on which the message arrived. A mismatch indicates a malfunctioning
    sender; the message shall be discarded.
    \item Apply the inline gossip to the local CRDT state as described in
    section~\ref{sec:session_crdt}.
    \item Deliver the application payload to every matching local subscriber,
    if any.
\end{enumerate}

Best-effort messages do not require receive-side deduplication at the session
layer, because the transport already deduplicates transfers and the publisher
never retransmits.

\section{Reliable message delivery}\label{sec:session_reliable}

A reliable publication is carried by a \texttt{msg\_rel} transfer and is
acknowledged by every intended recipient with a \texttt{msg\_ack}. The
publisher retransmits until every known recipient has acknowledged or the
deadline expires.

\subsection{Associations}

A publisher tracks its active recipients in an \emph{association set}. Each
association records:

\begin{description}
    \item[\textit{remote\_id}] The network-wide identifier of the remote node.
    \item[\textit{unicast\_ctx}] The transport-specific metadata needed to send
    a unicast message back to the remote, as supplied by the transport layer
    with each incoming acknowledgement.
    \item[\textit{last\_seen}] Timestamp of the most recent acknowledgement
    received from the remote.
    \item[\textit{slack}] A miss counter incremented each time the remote
    fails to acknowledge a reliable message within the deadline. The
    association is forgotten when \textit{slack} exceeds a small limit
    (recommended default: $2$).
\end{description}

Associations are created opportunistically on receipt of a positive
acknowledgement for a reliable message, and are removed when they age out of
the session lifetime (section~\ref{sec:session_parameters}) or when their
slack counter exceeds the limit.

\subsection{First publication}

When a publisher sends its first reliable message on a topic it may have no
associations at all. The first message is multicast on the topic's subject-ID
and the publisher treats the arrival of \emph{any} positive acknowledgement as
success; further acknowledgements arriving after the first continue to
populate the association set in the background.

\subsection{Retransmission policy}

The publisher retransmits a reliable message using an exponential backoff
schedule. The first retransmission is scheduled one \emph{ACK baseline
timeout} after the initial attempt; each subsequent retransmission doubles the
timeout. The final attempt is capped by the publication deadline supplied with
the publication. Implementations may scale the ACK baseline timeout by the
message's priority level, for example to give higher-priority publications
tighter timeouts.

Initial attempts are always multicast, since the publisher does not yet know
which recipients are live. As acknowledgements arrive and unacknowledged
recipients remain, the publisher may elect to unicast the remaining
retransmissions directly to the holdouts using the transport unicast metadata
recorded in the association entries. This multicast-to-unicast fallback is a
heuristic chosen to reduce network and subscriber load; it is not required for
interoperability.

\subsection{Acknowledgements}

A \texttt{msg\_ack} is sent by any subscriber that accepts a reliable message
for processing. The acknowledgement is unicast back to the publisher at the
priority of the original message, and references the message by its topic
hash and tag.

A \texttt{msg\_nack} is sent only in response to a reliable \emph{unicast}
message for which the receiver has no matching subscriber. It tells the
publisher to prune the association promptly so that further retransmissions
are not wasted on a recipient that cannot accept the topic. \emph{Multicast}
reliable messages do not generate NACKs, because subscription membership is
implicit in the transport's multicast join state: a subscriber that is no
longer interested simply stops participating in the multicast group.

\begin{remark}
    The asymmetry between unicast and multicast NACKs is intentional. For
    multicast delivery, the publisher learns that a remote has left by
    observing its absence from future acknowledgements; for unicast delivery,
    no such implicit signal exists, so an explicit NACK is required.
\end{remark}

\subsection{Receive-side deduplication}

A reliable message may be delivered more than once by the transport if an
acknowledgement is lost and the publisher retransmits the same tag. The
receiver eliminates such duplicates by keeping a per-remote cache keyed by
\texttt{(remote\_id, tag)} and discarding any tag that has already been
accepted. Dedup state is soft and is aged out after the session lifetime
(section~\ref{sec:session_parameters}).

\subsection{Observable states of a reliable publication}

A reliable publication transitions through the following observable states,
which are named here to describe protocol behaviour rather than any API:

\begin{description}
    \item[\textsc{pending}] The initial state; set when the first attempt is
    submitted to the transport. The publication remains pending until a
    terminal event occurs.
    \item[\textsc{done, success}] Reached when at least one positive
    acknowledgement has been received and all remaining associations have
    either acknowledged or aged out of the publication's deadline.
    \item[\textsc{done, delivery failure}] Reached when the publication
    deadline expires without any positive acknowledgement.
    \item[\textsc{done, explicit rejection}] Reached on receipt of a
    \texttt{msg\_nack} for a unicast reliable message.
    \item[\textsc{done, transient lag}] Reached when the publication
    deadline is about to expire but the publisher cannot determine whether
    the retransmission schedule actually delivered the message due to local
    scheduler pressure. This state is informational: the application is
    warned that reachability information for this publication is
    incomplete.
\end{description}

\section{Deduplication and ordering}\label{sec:session_dedup}

Session-layer deduplication applies to reliable messages only, as described in
section~\ref{sec:session_reliable}. Best-effort messages are deduplicated by
the transport and no additional session-layer work is required.

A subscriber may optionally enable \emph{per-remote ordered delivery}.
When enabled, the subscriber uses the monotonic per-publisher \texttt{tag}
field to reconstruct the publication order within a bounded reorder window.
Messages whose tag lies within the window are delivered in increasing tag
order; messages whose tag lies outside the window because they arrived too
late are dropped. Ordering is guaranteed per remote publisher only, not across
remotes, because there is no global clock.
